% =============================================================================
% Online Appendix to Section 7 (the opportunity-security space): further results and caveats.
% Numbers: paper/results/joint.json (scripts/paper/fig_joint.R, final rerun of 2026-10-12).
% Sample: cells rankable on both pillars (config/rankability.yml).
% =============================================================================
\section{The opportunity-security space: further results}\label{oa:joint}

\paragraph{Sample.} Of the 3{,}075 cells scored by both pillars but not rankable on both, 607
belong to territories and 2{,}468 to sovereign economies whose imputed weight share exceeds one
half in Energy Security only (721), in Economic Opportunity only (1{,}006) or in both (741). The
number of rankable economies per pair ranges from 40 (Green Hydrogen midstream) to 175 (Nuclear
upstream). In this sample technology-stage fixed effects account for 41.9 per cent of the variance
of $\ES_{its}$ and 53.4 per cent of that of $\EO_{its}$. The six EU members that enter the allied
network only through the EU node are not counted as coalition members.

\paragraph{Co-movement by technology.} The within-pair correlation exceeds 0.5 in 19 of the 30
pairs and is negative only in Green Hydrogen midstream. Pooled on levels it is 0.48 and on pooled
within-pair ranks 0.57. The stage gradient holds within wind (0.37, 0.72 and 0.89 by stage),
nuclear (0.16, 0.49, 0.89), coal (0.35, 0.56, 0.79), gas (0.40, 0.54, 0.78) and oil (0.50, 0.55,
0.75). Solar (0.48, 0.78, 0.58) and electric vehicles (0.60, 0.94, 0.80) peak midstream; in
batteries the correlation falls along the chain (0.63, 0.61, 0.40); in green hydrogen the midstream
value ($-0.21$) lies below the upstream one (0.03). Medians by technology run from 0.03 (green
hydrogen) to 0.80 (electric vehicles). Upstream the correlation is positive in all ten pairs. In the
five mineral pairs with a trade basket on both sides (batteries 0.63, grid 0.60, electric vehicles
0.60, solar 0.48, wind 0.37), opportunity's Trade and security's Minerals Trade measure exports of
overlapping materials and correlate at 0.70 to 0.85. The fossil-fuel pairs follow (oil 0.50, gas
0.40, coal 0.35); nuclear (0.16) has half of its economies at one security value, and green
hydrogen (0.03), the only upstream pair without a trade basket, rests on platinum-group endowments
on the security side and on cost, demand, readiness and production on the opportunity side.

\paragraph{Shared components.} A component counts as shared when the other pillar contains a
component of the same category with identical observed values: every component of Consumption,
Energy Access, Investment and Energy Prices, two of the three security and two of the five
opportunity components of Trade, and five of the six components of Foreign Dependency in each
pillar; no Production component. In seven downstream pairs (batteries, coal, electric vehicles,
gas, green hydrogen, nuclear, oil) the security score is a reweighting of part of the opportunity
score, which is why Green Hydrogen (0.92) and Nuclear (0.89) downstream show some of the highest
correlations. The non-shared median correlation has an interquartile range of 0.20 to 0.40, a 45
per cent reduction from the full pillars. The level covariance is positive in all 30 pairs,
including Green Hydrogen midstream. Of the pooled within-pair covariance, the two cross terms of
equation~\eqref{eq:joint_cov} account for 24 per cent and the non-shared term for 16 per cent
(58 per cent upstream, nothing downstream); in the median pair the shared term accounts for 19 per
cent. The cross terms are not definitional but run through shared inputs.

The non-shared rank correlation is 0.37 midstream and 0.30 upstream. The pillars' distinctive trade
components (net imports and, upstream, minerals trade for security; revealed comparative advantage,
export size and feasibility for opportunity) all reward exporters of the stage's basket; removing
opportunity's trade components as well lowers the medians to 0.15 midstream and 0.02 upstream. In
the three downstream pairs where it can be computed (geothermal, solar, wind) it is $-0.04$.

\paragraph{Category correlations.} Consumption, Energy Access and Investment correlate perfectly
with their counterparts; Trade, sharing two components, 0.62; Production, from different sources
and transformations of the same physical quantity, 0.27; and Foreign Dependency, whose shared
components are missing for many economies and imputed under different rules, 0.04. The median
correlation between categories with different names is 0.14 over 47 category pairs. Technology
Demand and Cost Competitiveness have absolute correlations of at most 0.12 and 0.17 with every
security category (the most negative is Cost Competitiveness with Reserves, $-0.17$). Within
security, Production and Reserves average 0.79 across seven upstream pairs (0.72 to 0.91 in the six
critical-mineral pairs, where 62 to 85 per cent of rankable economies record zero on both); Nuclear
upstream has the lowest value (0.61), because many holders of uranium resources do not mine them.

\paragraph{Regions.} Laggards outnumber leaders because a tied cell at a median counts as low and
ties are common where a pillar rests on a zero-anchored component. Across pairs, the observed share
of opportunity without security tracks the copula-implied one (correlation 0.90). Off-diagonal
shares are 19.5 and 17.2 per cent upstream, 13.8 per cent (opportunity without security)
midstream, and 12.4 and 10.2 per cent downstream; by technology the leapfrog share runs from 9.5
per cent (electric vehicles) to 20.2 per cent (green hydrogen and nuclear). In Green Hydrogen
midstream 34 of 40 economies share the lowest security value, which is also the median; in Nuclear
upstream 87 of 175 share 0.172, the median, and 37 of the 49 cells in the leapfrog region are such
economies. Counting ties as high empties the region in Green Hydrogen midstream and lowers it to
6.9 per cent in Nuclear upstream. The next pairs, Green Hydrogen upstream (24 per cent), Coal
upstream (22) and Batteries downstream (21), have few ties and shares near the copula value.

\paragraph{The coalition.} The partners, many of them major suppliers, have 21.4 per cent of cells
with security but not opportunity. The median core member has 16.7 per cent of its pairs in the
leapfrog region, the median partner 6.9 per cent. Austria, Portugal, Luxembourg (11 of 30 each),
Ireland (10) and Croatia (9) lead; without the two tie pairs their counts are 9, 10, 10, 8 and 7.
Among partners the largest counts are Thailand and the Philippines (5 of 29), then Malaysia, Viet
Nam and Qatar (4). Thirteen members have no leapfrog pair: Canada, Japan, the United Kingdom,
Australia and Ukraine in the core, and T\"urkiye, Brazil, Argentina, the Democratic Republic of the
Congo, the Dominican Republic, Gabon, Namibia and Ghana among partners. Brazil, a partner on its
reserves alone, is the clearest case of security without opportunity. The pairs most often in the
leapfrog region across members are Gas and Oil upstream and Green Hydrogen midstream (14 members
each), Coal upstream (13), Solar downstream (12), Nuclear upstream (11) and Grid upstream (10);
without the two tie pairs the coalition share falls from 12.8 to 11.8 per cent. Members' rankable
pairs number 22 to 30 in the core and 9 to 29 among partners, so the shares of Costa Rica, Gabon,
Namibia (9 pairs), Guatemala (10), the Dominican Republic and Ghana (11) rest on few cells.

\paragraph{Within countries.} Of the 174 countries with a Strategic Index, seven with fewer than
five cells are left out; the remaining 167 have 5 to 30 cells (median 22; 50 have fewer than ten).
The correlation is positive for 164 countries; the three exceptions (Afghanistan, Bhutan, the
Maldives) are not members. Among members it averages 0.52, from 0.27 (New Zealand) to 0.77 (Viet
Nam). Security varies more than opportunity across a country's sectors (mean standard deviations
0.15 and 0.07); the covariance share ranges from 16.5 per cent (Ghana) to 66 per cent (India) among
members, and the standard deviation of $\Delta_{its}$ from 0.068 (Lao PDR) to 0.26 (the Gambia).
The pairs' mean pillar values are only weakly related across pairs (0.33); after subtracting pair
means the within-country correlation rises to 0.67 and the standard deviation of the gap falls to
0.08, so part of the gap's spread is common to every country.

\paragraph{The Strategic Index and the regions.} Every Strategic Index cell (3{,}154 for 174
countries) lies in this sample. A country's top sector, before eligibility rules, lies in the
leapfrog region for 47 of 174 countries (27.0 per cent; 15 members), against 15.2 per cent of
cells. Nuclear upstream is the top sector of 62 countries, 43 of them at the common security value,
which gives a large need whichever side of the opportunity median they fall. After the endowment
screen, 169 countries keep an eligible cell and for 35 of them (20.7 per cent) the top one lies in
the region. Of the 57 allied picks (43 midstream), the eight in the region are Austria, Belgium,
Italy, the Netherlands, Romania and Viet Nam in Batteries midstream and Indonesia and Slovakia in
Nuclear midstream; 21 are leaders, 14 have security without opportunity and 14 are laggards, while
61.0 per cent of eligible sectors are leaders. None of the eight Nuclear upstream picks is in the
region: all eight members hold or mine uranium. Costa Rica, the Dominican Republic and Guatemala have
four eligible sectors each (nuclear, oil, gas and coal midstream) and Namibia five (the same and
Nuclear upstream). Midstream nuclear security is built from Trade alone and compressed around its
median (quartiles 0.435, 0.482, 0.530), so their security ranks of 0.50 to 0.71 lie at or above it,
and the technology weight (nuclear 0.999, gas 0.422, coal 0.024, oil zero) favours nuclear.

\paragraph{Caveats.} Percentile ranks discard distances, and the median split turns small
differences into differences of region; where many economies tie at a median the convention
decides the region of the whole group (counting ties as high lowers the upstream leapfrog share
from 19.5 to 15.5 per cent and the coalition's from 12.8 to 10.8). Regions are relative to the
rankable economies of each pair, a set that leans towards economies with good statistical
coverage. The rank-based measure of mechanical overlap covers 23 pairs; the covariance decomposition
covers all 30 but concerns levels and is dominated by the pairs with the largest covariances.
Within-country statistics compare category bundles that differ across pairs. The endowment screen
is a rule for the allied picks, not part of either pillar or of the index. More generally, shared
components make the implied trade-offs of a composite hard to see in the aggregate
\citep{ravallion2012mashup}; inspecting the correlation structure among indicators
\citep{oecd2008handbook} is what reveals them here.
